\chapter{Declarations}
\label{sec:statements}

\paragraph{Data and licences.} Table~\ref{tab:licences} summarises the terms of every source. The Qur'an and hadith texts
are fetched from the public QuranEnc and HadeethEnc interfaces with their references and links preserved; their
redistribution terms are not stated in the files we hold and will be confirmed with the publishers before any release
of the corpora. Sign data are used for research. Poses derived from sources marked non-commercial, no-derivatives or
no-redistribution (ASL Citizen, Isharah, WSLP) are never published; public-video sources without a stated licence are
used only inside the prototype, and their owners will be asked for permission before any release. The YouTube-ASL
keypoints are CC BY 4.0, so their converted form is shared with attribution.

\begin{table}[ht]
\centering\small
\caption{Licence summary.}
\label{tab:licences}
\begin{tabularx}{\textwidth}{L{4.6cm} L{4.2cm} Y}
\toprule
Source & Licence & Use in Isharati \\
\midrule
YouTube-ASL Clip Keypoints \cite{youtubeaslkp} & CC BY 4.0 & converted, shared with attribution (1 of 10 parts so far) \\
How2Sign \cite{how2sign} & CC BY-NC 4.0 (dataset); MIT (README of the feature files used) & recogniser training (planned) \\
CISLR \cite{cislr} & AFL-3.0 & ISL lexicon \\
WLASL \cite{wlasl}, MS-ASL \cite{msasl} & C-UDA, research use & ASL lexicon \\
KArSL \cite{karsl} & research use & ArSL lexicon \\
ASL Citizen \cite{aslcitizen} & Microsoft Research licence, non-commercial, no redistribution & ASL lexicon, not redistributed \\
WSLP \cite{wslp} & CC BY-NC-ND 4.0 & ISL lexicon, poses not published \\
Isharah \cite{isharah} & CC BY-NC-ND 4.0 & experiments only, not in the lexicon \\
iSign \cite{isign} & CC BY-NC-SA 4.0 & noted for future use \\
ArabSign \cite{arabsign} & not stated in the paper & evaluation (gold glosses) and recogniser data \\
YouTube channels (GDM, Noor, Obscure ASL, Tawasol, Al-Kharj, Saudi dictionary, ArabicSignLanguage, DisabilityApps, Ayman Abbas,
Mrami, T\.{I}D S\"ozl\"u\u{g}\"u, TiDiSLaM) & not stated & prototype only; permission to be requested \\
VRM avatar (pixiv sample) \cite{vrm} & VRM licence 1.0: modification, redistribution, commercial and religious use allowed & default avatar \\
\bottomrule
\end{tabularx}
\end{table}

\paragraph{Data availability.} The code, the merged lexicon indexes (glosses, sources and review status, without the
restricted poses), the evaluation scripts and the result files are in the project repository (private during the challenge). The converted YouTube-ASL keypoints are being published
as a Hugging Face dataset with the original CC BY 4.0 attribution.

\paragraph{Ethics.} The system answers religious questions only from cited sources and refers ruling questions to
scholars. No personal data are collected by the application; questions are processed locally, except for the calls to
the LLM providers (NVIDIA NIM, Groq), which receive the question and the retrieved passages. Sign recordings are of
signers who published them; they are used for research, and poses are not redistributed where the licence forbids it.
Outputs are labelled as pending review until Deaf signers approve them.

\paragraph{Author contributions (CRediT).} \emph{Fatimah Emad Eldin}: conceptualization, methodology, software, data
curation, formal analysis, investigation, visualization, writing (original draft). \emph{Asmaa Al-Mirghani}:
conceptualization, validation of religious content and sources, methodology (evaluation design), writing (review and
editing).

\paragraph{Conflicts of interest.} The authors declare no conflict of interest.

\paragraph{Funding.} This work received no external funding.

\paragraph{Use of AI tools.} Large language models are components of the system (query expansion, answer writing under
the grounding rules, glossing, and proposals of concept mappings), as described above. AI coding and writing assistants
were used to develop the software and to draft this document; the authors reviewed the content, and every number was
taken from the project's files.
